\documentclass[]{article}

%opening
\title{}
\author{}

\begin{document}

\maketitle

\begin{abstract}

\[
\textbf{Soluciones}
\]

\begin{enumerate}
	% 1
	\item \textbf{Resuelva:}
	\[
	\frac{(x^{6}y^{3})^{-1/3}}{(x^{4}y^{2})^{-1/2}}
	= (x^{6})^{-1/3}(y^{3})^{-1/3}\,(x^{4})^{1/2}(y^{2})^{1/2}
	= x^{-2}y^{-1}\cdot x^{2}y
	= 1.
	\]
	\[
	\sqrt[5]{\frac{8x^{3}}{y^{4}}}\,\sqrt[5]{\frac{4x^{4}}{y^{2}}}
	= \sqrt[5]{\frac{32x^{7}}{y^{6}}}
	= \sqrt[5]{2^{5}}\;\sqrt[5]{\frac{x^{7}}{y^{6}}}
	= 2\cdot \frac{x}{y}\,\sqrt[5]{\frac{x^{2}}{y}}.
	\]
	
	% 2
	\item \textbf{Exprese como polinomio} \; $(2x-1)(x^2-5)(x^3-1)$:
	\[
	(x^2-5)(x^3-1)=x^5-5x^3-x^2+5,
	\]
	\[
	(2x-1)(x^5-5x^3-x^2+5)
	=2x^6-10x^4-2x^3+10x -x^5+5x^3+x^2-5,
	\]
	\[
	\boxed{\,2x^6-x^5-10x^4+3x^3+x^2+10x-5\,}.
	\]
	
	% 3
	\item \textbf{Simplifique} \; $\displaystyle \frac{2x+1}{x^2+4x+4}-\frac{6x}{x^2-4}+\frac{3}{x-2}$:
	\[
	x^2+4x+4=(x+2)^2,\qquad x^2-4=(x-2)(x+2).
	\]
	Con denominador común $(x-2)(x+2)^2$,
	\[
	\frac{(2x+1)(x-2)-6x(x+2)+3(x+2)^2}{(x-2)(x+2)^2}
	=\frac{-x^2-3x+10}{(x-2)(x+2)^2}.
	\]
	Como $-x^2-3x+10=-(x+5)(x-2)$, se obtiene
	\[
	\boxed{\,-\dfrac{x+5}{(x+2)^2}\,}\qquad (x\ne \pm 2).
	\]
	
	% 4
	\item \textbf{Resuelva} \; $\displaystyle \frac{2}{2x+5}+\frac{3}{2x-5}=\frac{10x+5}{4x^2-25}$:
	\[
	\frac{2(2x-5)+3(2x+5)}{(2x+5)(2x-5)}
	=\frac{10x+5}{(2x+5)(2x-5)}
	=\frac{10x+5}{4x^2-25}.
	\]
	La igualdad es identidad siempre que los denominadores no sean cero:
	\[
	\boxed{\,x\in\mathbb{R}\setminus\left\{-\tfrac{5}{2},\tfrac{5}{2}\right\}\,}.
	\]
	
	% 5
	\item \textbf{Aislamiento:} costo \$2400; ahorro mensual $0.10\cdot\$200=\$20$.
	\[
	\text{Meses}=\frac{2400}{20}=120\quad\Rightarrow\quad \boxed{120\ \text{meses (10 años)}}.
	\]
	
	% 6
	\item \textbf{Millas en ciudad:} Sea $c$ ciudad y $h$ carretera.
	\[
	\begin{cases}
		c+h=1800,\\[2pt]
		\dfrac{c}{25}+\dfrac{h}{40}=51.
	\end{cases}
	\Rightarrow \frac{1800-h}{25}+\frac{h}{40}=51
	\Rightarrow 14400-3h=10200
	\Rightarrow h=1400,
	\]
	\[
	\boxed{c=1800-1400=400\ \text{millas}}.
	\]
	
	% 7
	\item \textbf{Escriba en la forma $a+bi$:}
	\[
	(3+4i)(3-4i)=3^2+4^2=\boxed{25}.
	\]
	\[
	\frac{-4+6i}{2+7i}\cdot\frac{2-7i}{2-7i}
	=\frac{(-4+6i)(2-7i)}{2^2+7^2}
	=\frac{34+40i}{53}
	=\boxed{\frac{34}{53}+\frac{40}{53}\,i}.
	\]
	
	% 8
	\item \textbf{Perro a más de 9 pies:} $s=16t^2+24t+1>9$.
	\[
	16t^2+24t-8>0 \;\Longleftrightarrow\; 2t^2+3t-1>0.
	\]
	Raíces: $t=\dfrac{-3\pm\sqrt{9+8}}{4}=\dfrac{-3\pm\sqrt{17}}{4}$.
	Como la parábola abre hacia arriba, 
	\[
	\boxed{\,t>\dfrac{-3+\sqrt{17}}{4}\approx 0.281\ \text{s}\,}
	\quad (\text{para } t\ge 0).
	\]
	
	% 9
	\item \textbf{Prisma recto con base equilátera $a=6$ cm y altura $H=12$ cm:}
	\[
	A_{\triangle}=\frac{\sqrt{3}}{4}a^2=\frac{\sqrt{3}}{4}\cdot 36=9\sqrt{3}\ \text{cm}^2,
	\]
	\[
	V=A_{\triangle}H=9\sqrt{3}\cdot 12=\boxed{108\sqrt{3}\ \text{cm}^3},
	\]
	\[
	A_L=(\text{perímetro})\cdot H=(3a)H=18\cdot 12=216,
	\quad A_T=2A_{\triangle}+A_L=18\sqrt{3}+216.
	\]
	\[
	\boxed{A_{\text{total}}=216+18\sqrt{3}\ \text{cm}^2}.
	\]
	
	% 10
	% 10
		\item \textbf{Área sombreada:} con la información dada (un segmento oblicuo que forma $30^\circ$ con la base,
	un triángulo rectángulo grande y otro pequeño, y la longitud total $4$ en la base)
\[
\text{Dado: } b=4,\qquad \theta=30^\circ.
\]
La altura correspondiente al ángulo de \(30^\circ\) es
\[
h=b\cdot \tan(30^\circ)=4\cdot \tan(30^\circ)=4\cdot\frac{1}{\sqrt{3}}=\frac{4}{\sqrt{3}}.
\]
El área del triángulo sombreado es
\[
A=\frac{1}{2}bh=\frac{1}{2}\cdot 4\cdot \frac{4}{\sqrt{3}}
=\frac{8}{\sqrt{3}}.
\]
Si se racionaliza el denominador:
\[
A=\frac{8}{\sqrt{3}}\cdot\frac{\sqrt{3}}{\sqrt{3}}
=\frac{8\sqrt{3}}{3}\ \text{(unidades}^2\text{)}.
\]
Aproximando numéricamente:
\[
A\approx \frac{8}{1.732}=4.619\ (\text{unidades}^2).
\]



\[
\textbf{Datos:}\qquad AB=4,\quad \angle B=30^\circ,\quad \angle C=90^\circ.
\]

\text{En un triángulo rectángulo }30^\circ\text{-}60^\circ\text{-}90^\circ\text{ con hipotenusa }h:
\]
\[
\text{cateto opuesto a }30^\circ=\frac{h}{2},\qquad
\text{cateto opuesto a }60^\circ=\frac{h\sqrt{3}}{2}.
\]

\text{Aquí }h=AB=4 \Rightarrow
\]
\[
AC=\frac{4}{2}=2,\qquad BC=\frac{4\sqrt{3}}{2}=2\sqrt{3}.
\]

\text{Coordenadas convenientes (para claridad):}
\[
A=(0,0),\quad B=(4,0).
\]
Si \(C=(x_C,y_C)\), de las distancias
\[
AC^2=x_C^2+y_C^2=2^2=4,\qquad BC^2=(4-x_C)^2+y_C^2=(2\sqrt{3})^2=12,
\]
restando se obtiene \(x_C=1\) y por tanto \(y_C=\sqrt{3}\). Así
\[
C=(1,\sqrt{3}).
\]

\bigskip
\textbf{Triángulo sombreado grande (lado derecho).}  
La recta que pasa por \(A\) y \(C\) tiene pendiente \(m=\dfrac{y_C-0}{x_C-0}=\dfrac{\sqrt{3}}{1}=\sqrt{3}\).
Como esa recta es perpendicular a la recta \(CB\) (se aprecia en la figura)
y corta la recta vertical que pasa por \(B\) (es decir la recta \(x=4\)),
el vértice superior del triángulo sombreado derecho es
\[
T=(4,\; y_T),\qquad y_T = y_C + m(4-x_C)=\sqrt{3}+\sqrt{3}\cdot(4-1)=4\sqrt{3}.
\]

El triángulo sombreado derecho tiene vértices \(T=(4,4\sqrt{3}),\; B=(4,0),\; C=(1,\sqrt{3})\).
Como \(TB\) es vertical (longitud \(4\sqrt{3}\)) y la distancia horizontal de \(C\) a la recta \(x=4\) es \(4-1=3\),
su área es
\[
A_{\text{sombra derecha}}=\frac{1}{2}\cdot \text{(altura }TB)\cdot \text{(distancia horizontal)} 
= \frac{1}{2}\cdot 4\sqrt{3}\cdot 3 = 6\sqrt{3}.
\]

\bigskip
\textbf{Triángulo sombreado izquierdo.}  
En la construcción geométrica exacta la recta que contiene a \(A\) y \(C\) pasa por \(A\) (es decir \(A,C,T\) son colineales), por lo que la «pequeña» región sombreada a la izquierda es degenerada en el modelo geométrico ideal (su área es 0). La pequeña área visible en la imagen es sólo un efecto del grosor del trazo del dibujo.

\bigskip
\textbf{Respuesta:}
\[
\boxed{\,A_{\text{sombreada total}}=A_{\text{sombra derecha}}=6\sqrt{3}\ \text{(unidades}^2\text{)}\approx 10.3923.}
\]


\end{enumerate}


\end{abstract}

\section{ esto es una seccio }

\end{document}
